% Source: https://artofproblemsolving.com/community/c3h2671792_500th_post
% Stable identity: aops-topic-2671792-q06
In each triangle the angle at $O$ is $30^\circ$ and the hypotenuse is $\dfrac{2}{\sqrt3}$ times the longer leg. The first triangle has $OA=\sqrt3$ and $OB_1=2$, and in general
\[OB_k=\sqrt3\left(\dfrac{2}{\sqrt3}\right)^k.\]
The ray $OB_k$ makes an angle of $30k^\circ$ with $OA$, so the hypotenuse lies along $OA$ for the first time when $k=12$ (a full turn). Then
\[OB_{12}=\sqrt3\cdot\dfrac{2^{12}}{3^6}=\dfrac{4096\sqrt3}{729}>\sqrt3=OA,\]
so $OB_{12}$ contains $OA$, and the part not on $OA$ is the segment $AB_{12}$ (red in the sketch) of length
\[OB_{12}-OA=\dfrac{4096\sqrt3}{729}-\sqrt3=\boxed{\dfrac{3367\sqrt3}{729}}\approx8.0.\]

\begin{asy}
usepackage("amsmath");
size(6.5cm);
defaultpen(fontsize(10pt));
real s=sqrt(3);
pair Q(int k){return s*(2/s)^k*dir(30*k);}
pair O=(0,0);
for(int k=2;k<=11;++k) draw(O--Q(k-1)--Q(k)--cycle,gray(0.55));
draw(O--Q(0)--Q(1)--cycle,linewidth(0.9));
draw(O--Q(11)--Q(12)--cycle,linewidth(0.9));
draw(Q(0)--Q(12),red+linewidth(1.6));
dot(O);dot(Q(0));dot(Q(12));
label("$O$",O,SW);label("$A$",Q(0),SE);label("$B_1$",Q(1),NE);label("$B_{12}$",Q(12),E);
label("$1$",(Q(0)+Q(1))/2,E,fontsize(8pt));label("$OA=\sqrt3$",(5.2,-1.3),fontsize(8pt));
\end{asy}
